\documentclass[a4paper,twoside]{article}

\usepackage{epsfig}
\usepackage{subcaption}
\usepackage{calc}
\usepackage{amssymb}
\usepackage{amstext}
\usepackage{amsmath}
\usepackage{amsthm}
\usepackage{multicol}
\usepackage{pslatex}
\usepackage{apalike}
\usepackage{url}
\usepackage{algorithm2e}
\usepackage[bottom]{footmisc}
\usepackage{SCITEPRESS}     
\usepackage{makecell} 
\usepackage{booktabs} 
\usepackage{multirow} 
\usepackage{array} 
\usepackage{stfloats}
\usepackage{float}

\begin{document}

\newcommand{\B}[2]{\ifnum#1=1 \textbf{#2}\else #2\fi}

\title{Robustness of Bayesian Neural Networks Design Choices Against Weight and Bias Perturbation Induced by Single Event Upsets (SEUs)}

\author{\authorname{Revalda Putawara\orcidAuthor{0009-0008-7938-185X}, Dalila O'Grady\orcidAuthor{0000-0002-7982-2917},  James Pope\orcidAuthor{0000-0003-2656-363X}}
\affiliation{Intelligence Systems Laboratory, University of Bristol, UK }
% \affiliation{\sup{2}Department of Computing, Main University, MySecondTown, MyCountry}
\email{ey24622@bristol.ac.uk, mv22726@bristol.ac.uk, james.pope@bristol.ac.uk}
}

\keywords{Bayesian Neural Networks, Robustness, Single Event Upsets, Prior Distributions}

%\abstract{
%Recent theoretical machine learning research has shown that the traditional U-shaped bias-variance trade-off hypothesis is not correct for certain deep learning models. Complex models with more parameters will fit the training data well, often with zero training loss, but generalise poorly, a situation known as overfitting. However, some deep learning models have shown to generalise even after overfitting, a situation known as the double descent phenomenon. It is important to understand which deep learning models exhibit this phenomenon for practitioners to design and train these models effectively. It is not known whether more recent deep learning models exhibit this phenomenon. In this study, we investigate double descent in three recent neural network architectures: Liquid Time-Constant Networks (LTCs), Quantised Neural Networks (QNNs), and Spiking Neural Networks (SNNs). We conducted experiments on the MNIST, Fashion MNIST, and CIFAR-10 datasets by varying the widths of the hidden layers while keeping other factors constant. Our results show that LTC models exhibit a subtle form of double descent, while QNN models demonstrate a pronounced double descent on CIFAR-10. However, the SNN models did not show a clear pattern. Interestingly, we found the learning rate scheduler, label noise, and training epochs can significantly affect the double descent phenomenon.}

%\abstract{
%Bayesian Neural Networks (BNNs) model uncertainty by treating weights and biases as probability distributions, making them attractive for deployment in environments (\textbf{where uncertainly needs to be modelled?}) where reliability under hardware faults is critical. A key fault source is Single Event Upsets (SEUs), caused by radiation that can occur in harsh environments. Prior work has studied deterministic CNNs under SEUs and BNNs under \textbf{adversarial attacks}, but little research examines BNN robustness to SEUs. The problem is that current neural network architectures are not robust.  Previous research has shown typical deep neural networks fail when presented with perturbed (noisy) inputs, particularly in adversarial settings.  Even less research has been conducted into model perturbations in harsh environments.  Making neural networks more robust remains a challenging and active research area.  As space exploration expands, artificial intelligence will increasingly be used in these harsh environments.  Our solution ... In this study, deterministic and Bayesian CNNs were trained with four design choices (baseline, modified pooling, dropout, weight decay) and seven activation functions. For BNNs, three priors (Gaussian, Laplace, Uniform) were tested. Models were subjected to SEU injections across layers, parameters, bit positions, and neuron locations. Robustness was quantified using a new Aggregate Robustness Index (ARIn), combining accuracy and softmax differences. Results show deterministic CNNs achieved higher baseline accuracy (97.55\% vs. 82.85\%), but BNNs were more robust (ARIn 0.0779 vs. 0.1909). Dropout yielded the largest robustness gain, with Gaussian priors outperforming Laplace and Uniform. SEUs on variance parameters were more harmful than on means, and fully connected layers proved most vulnerable. Sigmoid was the most robust activation despite an accuracy trade-off. Overall, BNNs provide more stable prediction confidence, making them suitable for safety-critical applications such as medical imaging and ensemble systems.
%}

\abstract{
Bayesian Neural Networks (BNNs) model uncertainty by treating weights and biases as probability distributions, making them attractive for deployment in environments where uncertainty can accommodate unexpected perturbations, leading to improved robustness. A key fault source is Single Event Upsets (SEUs), caused by radiation that can occur in harsh environments. Although robustness to input perturbations has been widely studied for both deterministic and Bayesian CNNs, weight and bias perturbations is not as widely explored, particularly for Bayesian CNN models. The growing use of AI in space missions and other fault-prone environments heightens the urgency in understanding the robustness of BNNs to weight/bias perturbations like SEUs. To address this gap, we offer the solution of introducing Bayesian Neural Networks in harsh environments by simulating it under SEU injections and exploring which design gives better robustness. In this study, deterministic and Bayesian CNN models were subjected to SEU injections across layers, parameters, bit positions, and neuron locations. Robustness was quantified using a new Aggregate Robustness Index (ARIn), combining accuracy and softmax differences. Results show deterministic CNNs achieved higher baseline accuracy (97.55\% vs. 82.85\%), but BNNs were more robust (ARIn 0.0779 vs. 0.1909). Dropout yielded the largest robustness gain, with Gaussian priors outperforming Laplace and Uniform. SEUs on variance parameters were more harmful than on means, and fully connected layers proved most vulnerable. Sigmoid was the most robust activation despite an accuracy trade-off. Overall, BNNs provide more stable prediction confidence, making them suitable for safety-critical applications such as medical imaging and ensemble systems.
}

%were trained with four design choices (baseline, modified pooling, dropout, weight decay) and seven activation functions. For BNNs, three priors (Gaussian, Laplace, Uniform) were tested.

%Prior work has studied robustness of deterministic CNNs under SEUs and BNNs under \textbf{adversarial attacks}, but little research examines BNN robustness to SEUs.

%The problem is that current neural network architectures are not robust. Previous research has shown typical deep neural networks fail when presented with perturbed (noisy) inputs, particularly in adversarial settings. Even less research has been conducted into model perturbations in harsh environments. Making neural networks more robust remains a challenging and active research area.  As space exploration expands, artificial intelligence will increasingly be used in these harsh environments. 

\onecolumn \maketitle \normalsize \setcounter{footnote}{0} \vfill

% \section{\uppercase{Introduction}}
% \label{sec:introduction}

\section{Introduction}
\label{sec:introduction}
%\subsection{Robustness of Machine Learning Systems Overview}
\noindent

%TODO

% ARE WE COVERING THESE IN THE INTRODUCTION? RESULTS?
% \begin{enumerate}
% \item Machine learning systems are also deployed in harsh environments.
% \item SEEs (including SEUs) can happen in harsh environments
% \item Robustness is a performance quality of machine learning systems, which consists of security and reliability
% \item Motivation on choosing Bayesian CNNs : Bayesian CNN offers uncertainty and already proven to be more robust against adversarial attacks
% \item Unfortunately, no significant research has thoroughly look at how robust it is against SEUs
% \item Thus, the research gap is explored in this research
% \end{enumerate}

Machine learning (ML) models are widely used in our daily life. Most of ML systems can be found close to us, like a recommender system, large language model or even autonomous vehicle that utilizes Convolutional Neural Network (CNN) models to steer clear from obstacles. While most of those systems are deployed in a normal environment, some of ML models are also deployed in harsh environments with its own specific designated purpose. 

%One of machine learning system that is deployed in a harsh environment is CloudScout \cite{Giuffrida2020-ry}. CloudScout, which is installed in Hyperscout-2 satellite, part of the European Space Agency (ESA) PhiSat-1 mission, is a first in-orbit demonstration of Deep Neural Network for data processing on edge. It uses CNN to classify hyperspectral satellite images as either cloudy or not cloudy directly on-board, and only transmit the non cloudy image to the ground, significantly reducing bandwidth consumption.

A representative example is CloudScout \cite{Giuffrida2020-ry}, a CNN-based system onboard the ESA PhiSat-1 mission that performs on-board cloud detection on hyperspectral images, transmitting only non-cloudy data to the ground, reducing bandwidth usage.

Machine learning systems deployed in harsh environments are exposed to risks during its inference, and a common term to associate this performance with Deep Neural Networks is \textit{robustness} \cite{shafique_2021}. Robustness are categorized into two subproperties, security and reliability. While the security aspect of DNNs focuses on preventing information inside the network to be stolen or protecting the system to get an incorrect input (adversarial attack), the reliability aspect of robustness focuses on keeping the model to display a consistent output, parameter, or behaviour, despite changes of environmental condition or fabrication.

% \subsection{Motivations} %answer question 2 : Why it is important, and why should people care

Being deployed in the space, a harsh environment, we can assume that machine learning systems like CloudScout is prone to risks that exists due to being deployed outside of earth surface. One of the risks that could hamper the performance of machine learning models that are deployed in harsh environments is single event effects (SEEs).

One example of SEE that can happen to machine learning models deployed in a harsh environment is Single Event Upsets (SEUs). SEUs are a temporary change of memory or control bit \cite{petersen_single_2011}. One phenomenon that shows the characteristic of SEUs is the incident of bitflips.

Machine learning with a good quality of robustness, or Robust ML Systems, can mitigate the impact of SEUs/bitflips, rendering the output of the model to be consistent. In this research, we are motivated to understand the impact of SEUs/bitflips to weight/bias perturbation in Bayesian Neural Networks. 

This research is built upon work done in \cite{dennispoperobustcnn} that reviewed CNN design choices to develop robust ML models. This research is also motivated by previous research that has proved that Bayesian Neural Networks are more robust against adversarial attacks (security aspect of robustness). Although there are already a significant number of studies that addressed those two aforementioned fields, \textbf{there has been little to no investigation into the robustness of Bayesian Neural Networks against weight/bias perturbation} (reliability aspect of robustness). This dissertation seeks to address that gap.


% \subsection{Our Approach}

In this research, we will show multiple Bayesian Neural Network design choice and compare its performance measure. Subsequently, we will simulate Single Event Upset (SEUs) with bitflips across all Bayesian Neural Network design alternatives and measure its performance measures afterwards. Robustness is then measured from the difference on the result of the model, before and after the model is perturbed with SEUs.

We comprehensively compared deterministic and Bayesian CNN models' performance across 4 model variations, 7 activation functions, and 3 prior distributions.  In this research, we also proposed a new robustness measure called the Aggregate Robustness Index (ARIn), which combines accuracy change and softmax deviation. Overall, this study show whether Bayesian CNN model can achieve better robustness compared to deterministic CNN model or not. It also shows which model design that is more robust.


% \subsection{Study Aims}


% To find a Bayesian Neural Network design choice that is robust against perturbation of weight/bias, we aim to complete several tasks in this research:
% \begin{enumerate}
% \item Train multiple deterministic and Bayesian Neural Networks with different design choices and note their performance
% \item Do multiple experiments of SEU injection on the different deterministic and Bayesian Neural Network design and note its performance
% \item Compare the robustness of each different design
% \item Validate whether Bayesian Neural Networks are more robust compared to deterministic Neural Networks
% \end{enumerate}


% \subsection{Study Contributions}

%MAKE BULLETS?

This study provides the evaluation of BNN robustness under SEUs, which addresses current research gaps in prior work that has focused primarily on CNN robustness under SEUs and Bayesian CNN robustness under adversarial attacks.  The paper makes the following contributions:

\begin{itemize}
    \item We demonstrate a range of Bayesian Neural Network design choices that are robust against perturbation of the model's weight and  biases.
    \item We provide detailed analyse of the effects of SEUs across different attack variables, including layer type, weight versus bias, mean versus variance parameters, and bit index along with its original value.
    \item Propose a novel measure of model robustness for SEUs.  We show that the measure is capable of detecting critical aspects of models under SEUs attacks.
    % \item We show that deterministic CNN designs are more robust (according to accuracy measure) to attacks while BNNs are robust according to another measure (soft max difference).
    
\end{itemize}

% 


% To find a Bayesian Neural Network design choice that is robust against perturbation of weight/bias, we aim to complete several tasks in this research:
% \begin{enumerate}
% \item Train multiple deterministic and Bayesian Neural Networks with different design choices and note their performance
% \item Do multiple experiment of SEU injection on the different deterministic and Bayesian Neural Network design and note its performance
% \item Compare the robustness of each different design
% \item Validate whether Bayesian Neural Networks are more robust compared to deterministic Neural Networks
% \end{enumerate}

\section{Background}
\label{sec:relatedwork}

\subsection{Robust Machine Learning Systems}
\noindent

\subsubsection{Robustness against Adversarial Attack}

\iffalse
\begin{enumerate}
\item \cite{Blaas2021} Blaas 2021 
\item augmented robustness and the result of that research \cite{essbai2024a}
\item SEBR Robust Bayesian Neural Networks by Spectral Expectation Bound
Regularization : 2021
\item SGLD Bayesian Learning via Stochastic Gradient Langevin Dynamic: 2011
\item Carbone's research : June 2020
\end{enumerate}
\fi

Before we take a look into robustness of neural networks against SEUs/bitflips, we will see what are the defenses against adversarial attack applied in recent research, especially for Bayesian Neural Networks.

%\textbf{Stochastic Gradient Langevin Dynamics (SGLD)} did not address adversarial attack directly, but focuses on how deterministic model is overconfident to point estimate of weights used. SGLD was proposed in \cite{sgld2011} by adding noise during training while behaving similar to the Stochastic Gradient Descent (SGD). The final model predictions come from averaging many sampled parameter values. The robustness that is addressed in this paper is more related to how the model does not overfit and provide uncertainty, not reliability-related robustness. 

%In \cite{carbone2020a}, the author observed how Bayesian Neural Network is robust against \textbf{gradient-based adversarial attacks} e.g. Fast Gradient Sign Method (FGSM) and Projected Gradient Descent (PGD). This work explains that BNNs naturally mitigate gradient-based attacks because of the Bayesian posterior averaging.

%On the other hand, work in \cite{Blaas2021} claims the contrary, showing that BNNs trained with with accurate posterior inference (e.g. Hamiltonian Monte Carlo, advanced VI) \textbf{do not automatically gain adversarial robustness} compared to standard deterministic neural networks. This work showed an alternative result from \cite{carbone2020a}, in which the work shows that bayesian averaging does not wash out adversarial directions as strongly as the theory predicts. One of the proof is the non-vanishing gradients on the HMC sampling and susceptibility to PGD and FGSM attacks.

%\textbf{Spectral Expectation Bound Regularization (SEBR)} is introduced in \cite{sebr2021}. SEBR is a regularization technique added during training to limit the Lipschitz constant of BNN layers, which in the end improve robustness. It works against gradient-based attacks like FGSM and PGD by penalizing the expected spectral norm of each weight matrix (mean + variance), constraining how much the model output can change for a small input perturbations. The penalty term in SEBR is shown in $L_{\text{total}} = L_{\text{ELBO}} + \lambda \cdot L_{\text{SEBR}}$, in which $\lambda$ is the trade-off factor for the SEBR loss.


%\textbf{Stochastic Gradient Langevin Dynamics (SGLD)} \cite{sgld2011} introduces noise in training to capture uncertainty, although it is not specifically designed for adversarial robustness.

BNNs have been reported to mitigate gradient-based input attacks, e.g., Fast Gradient Sign Method (FGSM) and Projected Gradient Descent (PGD), due to posterior averaging \cite{carbone2020a}. There are also methods that employ multi-task loss combining unsupervised feature scattering and supervised margin loss to generate diverse and stronger adversarial examples for BNN training \cite{CHEN2022156}. \textbf{Spectral Expectation Bound Regularization (SEBR)} \cite{sebr2021} also increases robustness against gradient-based attacks by regularizing the spectral norm of weights. Contradicting work includes \cite{Blaas2021}, which shows that a more accurate posterior inference (e.g., Hamiltonian Monte Carlo, advanced VI) does not guarantee robustness.



%\begin{equation}
%L_{\text{total}} = L_{\text{ELBO}} + \lambda \cdot L_{\text{SEBR}}
%\label{eq:sebrpenalty}
%\end{equation}

\subsubsection{Robustness against Weight Perturbation}

%The term Single Event Upset Induced Parameter Perturbation (SIPP) is introduced in \cite{yan_2020a}. This paper reiterates SEU effect in flipping bits in memory or registers, especially in SRAM where model parameters are stored. It also shows that single bit-flip can severely degrade DNN accuracy, dropping accuracy from 70\% to 1\% in the experiment on using ResNet56 on CIFAR-100. This paper proposes hardware-level protection schemes : Triple Modular Redundancy (TMR), Error-Correcting Codes (ECC), and Selective Protection on top of TMR and ECC.

Single Event Upset Induced Parameter Perturbation (SIPP) \cite{yan_2020a} shows that SEUs in memory or parameters can drastically degrade accuracy. Proposed defense in this work includes Triple Modular Redundancy (TMR), Error-Correcting Codes (ECC), and Selective Protection on both.

%The work in \cite{dennispoperobustcnn} take a different approach in developing robust machine learning systems against weight perturbation by focusing on model-intrinsic robustness instead of hardware fault tolerance. This paper proposes two new bounded activation functions, namely Weighted Gaussian (WG) and Rectified Weighted Gaussian (RWG) that limits the effect of flipped bits. This paper also introduce smartpooling that ignores unreasonable outlier activations above threshold. Finally, this paper also shown that dropout introduced during training helped maintain accuracy while a smaller effect is also shown by introducing weight decay during training.

Model-intrinsic robustness is explored in \cite{dennispoperobustcnn}, where Weighted Gaussian (WG) and Rectified Weighted Gaussian (RWG) activation functions, smart pooling, dropout, and weight decay show improvements in deterministic CNNs under SEUs.

While both of the works mentioned previously take an interest in the robustness of deterministic neural networks, especially variants of convolutional neural networks, against weight/bias perturbation, \textbf{no significant research is found on observing Bayesian Neural Networks (BNNs) robustness against weight perturbation yet}. This area of research is what this paper aims to look into.

\subsection{Single Event Upsets}
\noindent

Single Event Effects (SEEs), as described in \cite{petersen_single_2011}, happened through the action of a single ionizing particle as it penetrates sensitive nodes inside an electronic device, illustrated in Figure \ref{fig:ionizationtIllustration}. It has been linked to a real-world system faults, for example, the 2008 Qantas A330 altitude change incident is suspected to be caused by a possible SEE \cite{ATSB2009QF72}. Some variants of SEEs are Single Event Upsets (SEUs), Multibit Upsets (MBUs), and Single Event Burnouts (SEBs).

%On October 2008, a Qantas A330 flight from Singapore to Perth experienced a sudden and unexpected altitude change, resulting in 36 crew and passengers injured. The Australian Transport Safety Bureau (ATSB) considers SEE as a possible cause of the accident \cite{ATSB2009QF72}. Some variants of SEEs are Single Event Upsets (SEUs), Multibit Upsets (MBUs), and Single Event Burnouts (SEBs).

%The effect of SEEs can be as small as a tiny glitch inside a monitor to a substantial error that is life threatening
%\cite{airsafeqantas}
%\cite{petersen_single_2011}

%Illustrations of heavy ions passing through a semiconductor, causing SEUs, is shown in Figure \ref{fig:ionizationtIllustration} (a) and (b). In Figure \ref{fig:ionizationtIllustration} (a), the ion creates a dense ionization track that generates enough charge to trigger an SEU. 

%In Figure \ref{fig:ionizationtIllustration} (b), a proton causes a nuclear reaction inside the device, producing secondary heavy ions that deposit charge and may also result in an SEU.

%SEUs or bit-flips have been simulated either by injecting non modifiable 0s or 1s into the weights of CNN layers at the software level to assess reliability and classify prediction deviations \cite{bosio2019}, or by injecting transient single and double bit flips into GPU registers and instruction outputs at the architecture level to study application-level Silent Data Corruptions (SDCs) \cite{hari2017sassifi}. Both studies show that SEUs can lead to incorrect or unsafe outputs in Deep Neural Networks, potentially degrading their reliability in safety-critical applications.  This research will follow a similar procedure to the SEU injection method in \cite{hari2017sassifi}, in which we will inject SEU (bitflip) on the weight/bias parameter of the model \textbf{when it is doing inference on a new data}. 

SEUs have been simulated either at the software level, by forcing bit values in CNN weights \cite{bosio2019}, or at the hardware level, by injecting transient bit-flips into GPU registers \cite{hari2017sassifi}. Both approaches show that SEUs can cause incorrect outpurs in deep networks, reducing reliability in safety-critical tasks. This study follows \cite{hari2017sassifi}, introducing bit-flips to model weights and biases \textbf{when it is doing inference on a new data}. 

%An example of the impact of bitflips to stored weight/bias is shown in Table \ref{table:robbothbitloc2}.

\begin{figure}[h]
\caption{Ionization path due to direct passage of heavy iron (a), and due to proton reaction in a device (b) as shown in \cite{petersen_single_2011}}
\includegraphics[scale=0.60]{images/ionization_path.png}
\centering
\label{fig:ionizationtIllustration}
\end{figure}

The IEEE-754 standard (\cite{IEEE754-2019}) defines binary (32, 64, 128 bits) and decimal (64, 128 bits) floating-point (FP) formats. Modern CPUs and frameworks (e.g. PyTorch) support binary machine numbers (FP32, FP64) for efficiency despite limited precision. Since Pyro is built on PyTorch, this research uses FP32, making it important to understand floating-point representation for analyzing the effects of Single Event Upsets (SEUs) on weight and bias parameters stored in FP32.

%\cite{pytorchNumericalAccuracy}
%\cite{pyroIntroductionPyro}

\subsection{Bayesian Neural Networks}
\noindent

In this section, we will take a look at the concept of Bayesian Neural Networks. This research assumes that the reader is familiar with the basics of Neural Networks, including the concept of feed-forward and backpropagation. The concept of Bayesian Neural Networks (BNNs) originated in the early 1990s. \cite{Mackay1992a} introduced a Bayesian framework for neural learning by trating weights and biases as probabilistic variables, addressing uncertainty and overfitting. 

%The idea of Bayesian Neural Networks (BNNs) dates back to some research done in the early 1990s. One of the early prominent works in Bayesian Neural Networks is the work done in \cite{Mackay1992a}. In that research, although not explicitly mentioning Bayesian Neural Network (BNN), David J.C. Mackay proposes a quantitative and practical Bayesian framework for learning in feedforward networks. This research revolutionizes how to think about weight and biases by giving them a probabilistic interpretation within a learning model. This work criticized the backpropagation framework's limitations, such as arbitrary hyperparameter settings, high data demands, and a lack of uncertainty estimates.

Equation \ref{eq:evidence} - \ref{eq:posterior} express the core of Bayesian learning. The marginal likelihood (Equation \ref{eq:evidence}) integrates over all weights configurations, embodying Occam's razor to balance model fit and complexity. Bayes' rule (Equation \ref{eq:posterior}) updates prior beliefs with data, yielding weight distributions that capture model uncertainty.

%The Bayesian framework integrates over weight distributions (Equation \ref{eq:evidence} - \ref{eq:posterior}).

%To fix the issues on the above points, Mackay proposed a Bayesian way of doing backpropagation by the following:

%\begin{enumerate}
%%\item Objective setting of control parameters: The bayesian framework uses evidence maximization to set the model hyperparameters initialized with associated priors. The marginal likelihood penalizes overfitting with Occam's razor, selecting the best model complexity by maximizing Equation (\ref{eq:evidence}).
%\item Objective setting of control parameters: In the Bayesian framework, model hyperparameters are determined by maximizing the marginal likelihood, which integrates over the parameter space with respect to the chosen priors. It automatically balances model fit and complexity, ensuring that the selected model achieves the best trade-off without overfitting.
%\begin{equation} \label{eq:occamsrazor}
%P(\mathcal{D} \mid \mathcal{M}) = \int P(\mathcal{D} \mid \mathbf{w}, \mathcal{M}) \, P(\mathbf{w} \mid %\mathcal{M}) \, d\mathbf{w}
%\end{equation}
%\item Demand of data size: Bayesian Inference penalizes complexity automatically, as shown in Equation (\ref{eq:evidence}), by automatically trades off data fit vs model complexity without the need for a separate validation or test set to detect overfitting.
%\item Lack of quantified uncertainty: Bayesian training infers a distribution over weights shown in Equation (\ref{eq:posterior}).

%The effective number of well-defined parameters ($\gamma$), are those weights, given the data, that are actually being used by the model in a meaningful way. With the above equation, well-defined parameters can be spotted by seeing which parameters have less variance, showing the model's confidence in their value.
%\end{enumerate}





\begin{equation}
P(\mathcal{D} \mid \mathcal{M}) 
= \int P(\mathcal{D} \mid \mathbf{w}, \mathcal{M}) \,
      P(\mathbf{w} \mid \mathcal{M}) \, d\mathbf{w}
\label{eq:evidence}
\end{equation}

\begin{equation}
p(\mathbf{w} \mid \mathcal{D}) \propto 
   P(\mathcal{D} \mid \mathbf{w}) \, p(\mathbf{w})
\label{eq:posterior}
\end{equation}


In \cite{Mackay1992a}, approximations are made to the posterior by using Laplace approximation and use Hessian of the loss to estimate posterior curvature. The term Bayesian Neural Network first coined in \cite{Neal1996bayesianlearning}. In that work, Radford M. Neal take the posterior distribution approximation to a new direction, by using Markov Chain Monte Carlo (MCMC) to not just do evidence maximization, but also sampling from the full posterior over weights, to approximate $p(\mathbf{w} \mid \mathcal{D})$. Illustration on how Bayesian Neural Networks differ to the Deterministic Neural Networks, in which each weight is assigned a distribution, is shown in Figure \ref{fig:BNNIllustration}.

\begin{figure}[h]
\caption{Illustration of Bayesian Neural Networks, as shown in \cite{blundell2015bayesbybackprop}}
\includegraphics[scale=0.30]{images/bayesbybackprop3.png}
\centering
\label{fig:BNNIllustration}
\end{figure}



\section{Methodology}
\label{sec:methodology}

%\section{\uppercase{Methodology}}

\subsection{Dataset}
\label{subs:dataset}

In this research, we are using an image dataset called "Ships in Satellite Imagery" \cite{robert_hammell_2018}, or "Shipsnet" in short. The dataset is satellite imagery captured from a commercial company called Planet. Those satellite images are taken from over the San Francisco Bay and San Pedro Bay areas of California. Initially, the pictures are orthorectified 3-meter pixel size, but for this task in detecting ships, 4000 80 x 80 RGB images are derived from it, labeled with either a "ship" (1) or "no-ship" (0) classification.

\begin{figure}[H]
    \centering
    %\begin{minipage}[b]{0.4\textwidth}
    %    \centering
    %    \includegraphics[width=\textwidth]{figure2/6.png}
    %    % (a)\\ % Add label (a) here
    %\end{minipage}
    %\hfill
    \begin{minipage}[b]{0.4\textwidth}
        \centering
        \includegraphics[width=\textwidth]{images/shipsnet_samples.png}
        % (b)\\ % Add label (b) here
    \end{minipage}
    \caption{Dataset}
    \label{fig:datasetsamples}
\end{figure}

%TODO explain about the preprocessing done for this dataset before it is deemed fit to be processed

The image is resized into 64 x 64 pixels, its pixel values are scaled, and then standardized, so that its value ranges from -1 to 1 and is centered around 0.


\subsection{Model architecture}

\textbf{Baseline.} The deterministic CNN model that we will set as a benchmark for this research is designed for binary image classification, as the class distribution dictated in subsection \ref{subs:dataset}. We will utilize 2 convolutional layers (namely \texttt{conv1} and \texttt{conv2}) with each followed by designated activation function and maxpooling (torch.nn.MaxPool2d). The output then flattened to be processed in a fully connected layer that outputs two class probability. In total, we will have 52,162 trainable weights and biases.

\textbf{Training.} The model is trained on Shipsnet dataset with 80/20 train-test split. The model will is trained on the dataset of batch size 16 in 20 epochs, keeping the model with the best test accuracy. Adam is used to minimize the loss function with learning rate of $10^{-3}$.

\textbf{Bayesian extension.} The Bayesian CNN model will follow the architecture of the deterministic CNN model, with a few differences. First, in this model, all weights and biases are modelled as random variable drawn from a prior distribution, which either Gaussian, Laplace, or Uniform prior. This variation will allow us to study how the tail behavior and spread of each prior influence robustness.

\textbf{Bayesian model training.} Our Bayesian CNN model will use negative Evidence Lower Bound (ELBO) loss. Given the dataset $D$ with observed inputs and labels, the ELBO loss aims to maximize the expected log-likelihood of the data from the sampled network weights while minimizing the KL divergence between the learned variational distribution $q(W)$ and the prior $p(W)$, trying to balance the data fit and regularization. ELBO is notated as follows:
\begin{equation}
\textbf{ELBO} \;\equiv\; 
\mathbb{E}_{q_\phi(W)}\big[\log p(D \mid W)\big]
\;-\;
\text{KL}\big(q_\phi(W) \;\|\; p(W)\big)
\label{eq:elbo_bnn}
\end{equation}

Bayesian Inference of these weights and biases during training will utilize Stochastic Variational Inference (SVI), with mean-field assumption. All weights and biases are assumed to be independent.

\textbf{Inference procedure.} Inference in Bayesian CNN models in this research is different from what is usually done in deterministic CNN models. We will use Monte Carlo Sampling in doing inference, with number of Samples $S=10$. Each iteration samples new weights ($W^{(s)}$) from a variational posterior family. Logits are averaged before applying argmax, reducing prediction variance. This mechanism shows the aggregation of Bayesian predictive uncertainty, which can not be found in the deterministic CNN. 

\subsection{Model Variations} \label{subs:modelvar}

%\subsubsection{Model variations}
%TODO

In this study, we will utilize 7 different activation functions to compare their robustness. Some of these are widely known activation functions like ReLu, Sigmoid, Tanh, and Sine. We also want to test the Weighted Gaussian (WG) and Rectified Weighted Gaussian introduced in \cite{dennispoperobustcnn} in this experiment, which are shown in Equation (\ref{eq:wg}) and Equation (\ref{eq:rwg}). As previous research shows that ReLU is the least robust activation function \cite{dennispoperobustcnn}, we also want to look at a different version of it called ReLU6, that clips ReLU to have a maximum value of 6.

\begin{minipage}{0.4\linewidth}
\begin{equation}
\text{WG}(x) = x e^{-x^{2}}
\label{eq:wg}
\end{equation}
\end{minipage}\hfill
\begin{minipage}{0.4\linewidth}
\begin{equation}
\text{RWG}(x) = \max(0, x e^{-x^{2}})
\label{eq:rwg}
\end{equation}
\end{minipage}

\begin{figure}[h]
\caption{Activation Functions Plot}
\includegraphics[scale=0.4]{images/activation_functions_new_stacked.png}
\centering
\label{fig:activationfunctions}
\end{figure}

To observe different model architecture design and their impact to the robustness of our model, we will explore a few options that we can choose to apply to that model, on top of our base model. The different model architecture designs are coded as follows:
\begin{enumerate}
\item Model 0: The base model.
\item Model 1: Instead of using maxpooling, the convolutional layer will be followed by smartpooling. This method is inspired from the work in \cite{santos_2019a}, that demonstrates that 16\% robustness is achieved in ReLU-based deterministic CNNs equipped with SmartPooling. In Smartpooling, if a value exceeds a threshold $\theta=10x$, the value will be discarded and the second highest value is chosen.
\item Model 2: Dropout layer is applied after conv2 layer during training. Dropout layer can benefit from mitigating the SEU-induced performance drop \cite{dennispoperobustcnn} by deactivating random neurons during training, with $p=0$ (off state) and $p=0.5$ (on state) in this experiment.
\item Model 3: Weight decay is implemented to the \texttt{Adam} optimizer with $\lambda = 10^{-4}$. Weight decay is a regularization technique that allows complex models to be trained on data sets of limited size without severe over-fitting \cite{bishoppatternrecognition2007}.
\end{enumerate}

While the deterministic CNN model have 28 model variations (7 activation functions and 4 architecture designs), the Bayesian CNN model that relies on sampling weight and bias from prior distribution will also be varied by the 3 different prior distribution (Gaussian, Laplace, and Uniform). This makes the Bayesian CNN model to have 84 model variations.

\subsection{Simulated SEU Process}

%TODO explain that the simulated SEU process will hit a bit index and change its value from 0 to 1 or vice versa.

In this research, we are using Floating Point 32 (FP32) to represent numbers in our machine. SEU injection will do a bitflip in one of the bit in FP32, changing the value from 0 to 1 or vice versa. Illustration of an SEU injection in a deterministic CNN model is as follows:

% added at 14/10/2025

\iffalse
\begin{itemize}
    \item Layers are denoted as \( m \in \{\mathrm{conv1}, \mathrm{conv2}, \mathrm{fc1}\} \).
        \item Each layer \( m \) has \( N_m \) output neurons.
    \item The neuron position under attack is denoted as \( \mathrm{pos} \in \{\mathrm{f}, \ell\} \),
    where \( \mathrm{f} \) corresponds to the first neuron and \( \ell \) to the last neuron, defined as:
    \[
    n_m^{\mathrm{f}} = 1, \qquad n_m^{\ell} = N_m.
    \]
    \item The parameter type is \( p \in \{\mathrm{w}, \mathrm{b}\} \) for weight or bias.
    \item The scalar value of weight or bias is represented as \( w_{m,n,p} \).
    \item The set of possible attacked bit indices follows the IEEE-754 32-bit floating-point representation. In this research, we limit the attack to these indices:
    \[
    I = \{0, 1, 3, 6, 10, 15, 21\}.
    \]
    %\item The binary representation of a 32-bit floating-point value \( w_{m,n,p} \) is denoted as
    %\[
    %\mathrm{bits}(w_{m,n,p}) = (b_{31}, b_{30}, \ldots, b_0),
    %\ b_i \in \{0, 1\}.
    %\]
    %The function \( \mathrm{FP32}(\cdot) \) converts a bit vector back to its floating-point scalar value.
\end{itemize}
\fi

%$m \ $

\begin{itemize}
    \item Layers are denoted as \( m \) and each layer has \( N_m \) output neurons.
    \item The neuron position under attack is denoted as \( \mathrm{pos} \in \{\mathrm{f}, \ell\} \),
    where \( \mathrm{f} \) corresponds to the first neuron and \( \ell \) to the last neuron, defined as:
    \[
    n_m^{\mathrm{f}} = 1, \qquad n_m^{\ell} = N_m.
    \]
    \item Parameter type is \( p \in \{\mathrm{w}, \mathrm{b}\} \) for weight or bias.
    \item The scalar value of weight or bias is represented as \( w_{m,n,p} \).
    \item A single attacked bit index \( i \) follows the IEEE-754 32-bit floating-point representation, \( i \in I \) 
    %In this research, we limit the attack to these indices:
    %\[
    %I = \{0, 1, 3, 6, 10, 15, 21\}.
    %\]
    %\item The binary representation of a 32-bit floating-point value \( w_{m,n,p} \) is denoted as
    %\[
    %\mathrm{bits}(w_{m,n,p}) = (b_{31}, b_{30}, \ldots, b_0),
    %\ b_i \in \{0, 1\}.
    %\]
    %The function \( \mathrm{FP32}(\cdot) \) converts a bit vector back to its floating-point scalar value.
\end{itemize}
%\vspace{1em}
For a selected layer \( m \), parameter type \( p \), neuron position \( \mathrm{pos} \),and bit index \( i \), the SEU targets the scalar parameter belonging to neuron \( n = n_m^{\mathrm{pos}} \).
Its bit representation is:
\[
\mathrm{bits}\big(w_{m,n,p}\big) = (b_{0}, \ldots, b_i, \ldots, b_{31})
,\ b_i \in \{0, 1\}.
\]

%Let \( j \in \{0, 1, \ldots, 31\} \) denote all the bit position index, the 

Let bit \( j \in \{0, 1, \ldots, 31\} \) be all bit value after bit \( i \) is attacked by SEU, according to the following rule:

\begin{equation}
b'_j =
\begin{cases}
1, & j = i \text{ and } b_i = 0, \\[6pt]
0, & j = i \text{ and } b_i = 1, \\[6pt]
b_j, & j \neq i.
\end{cases}
\label{eq:bitflip_case}
\end{equation}

The perturbed (post-SEU) scalar value $\tilde{w}_{m,n,p}$ is then defined as FP32 representation of its bits:

\begin{equation}
%\tilde{w}_{m,n,p}^{(i,\mathrm{pos})}
\tilde{w}_{m,n,p}
= \mathrm{FP32}\big(b'_{0}, b'_{1}, \ldots, b'_{31}\big)
%,\qquad n = n_m^{\mathrm{pos}}, \; i \in I.
\label{eq:perturbed_weight}
\end{equation}

For our deterministic CNN models, we will do 2,352 SEU injections. This SEU injections are the combination of 28 different models that we have (4 model variations and 7 activation functions), with each model get 84 different SEU injections, where each attack is specified by \( m \in \{\mathrm{conv1}, \mathrm{conv2}, \mathrm{fc1}\} \), \( \mathrm{pos} \in \{\mathrm{f}, \ell\} \), \( p \in \{\mathrm{w}, \mathrm{b}\} \), and \(i \in I = \{0, 1, 3, 6, 10, 15, 21\} \).

%which comprises a combination of 2 neuron attack locations (first/last neuron), 3 layer (conv1/conv2/fc1), 2 layer parameter (weight/bias), and 7 bit index. 

In our Bayesian CNN models, each parameter \( w_{m,n,p} \) is treated as a random variable
drawn from a variational posterior distribution. Let's take an example for a draw in a Gaussian distribution for Monte Carlo sample \(s \in S\):
\[
w_{m,n,p}^{(s)} \sim q_{\phi}(w_{m,n,p})
= \mathcal{N}\!\left(\mu_{m,n,p},\, \sigma_{m,n,p}^{2}\right).
\]
An SEU attack is applied before the weight/bias is sampled. Instead of applying SEU to \( w_{m,n,p}^{(s)} \), it is applied to either \( \mu_{m,n,p} \) (prior center parameter) or \( \sigma_{m,n,p}^{2} \) (prior spread parameter). The SEU attack procedure follows the procedure that has been explained earlier in this subsection.

For our Bayesian CNN models, we will do 15,456 SEU injections. 84 different models will be attacked with 84 ($7 \times 12$) attacks for the prior center and 72 ($6 \times 12$) attacks for the prior spread, for a total of 13,104 SEU injections. The combination of attacks is similar to what we have in the deterministic CNN models. The only difference is for the prior spread parameter, the bit index attacked is limited to \(i \in I = \{1, 3, 6, 10, 15, 21\} \). This is due to hitting the bit index 0 for prior spread will change the prior spread parameter sign from + to -, which is not valid to represent the spread parameter of our prior.

%\( w_{m,n,p}^{(s)} \),resulting the perturbed value \( \tilde{w}_{m,n,p}^{(s)} \)defined in Equation~\ref{eq:perturbed_weight}.

%For Bayesian CNN models, 84 different models will be attacked with 156 different SEU injections, 84 ($7 \times 12$) for the prior center and 72 (6 \times 12) for the prior spread, for a total of 13,104 SEU injections. The attack comprises of 36 different attacks, which combination comprises of 2 neuron attack locations, 3 layers and 2 layer parameters. We can only attack 6 variations of bit index into the prior spread, as hitting the bit index 0 for prior spread will change the prior spread parameter sign from + to -, which is not valid to represent the spread parameter of our prior.

The total SEU injection attacks that we will do in this study is 15,456 injections.

\subsection{Robustness Measure}

To evaluate how well a neural network withstands adversarial attack or weight/bias perturbation, we need a quantitative metric that can measure its robustness. This measure will help us in comparing the variation of models or strategies with each other and decide which ones are the most robust against perturbations.

The first metric that we can use to measure robustness is accuracy delta/accuracy difference. It is the most commonly used measure of robustness in this field of research and can be found in \cite{dennispoperobustcnn}, \cite{feng2024attackingbayesadversarialrobustness}, and \cite{pang2021evaluatingrobustnessbayesianneural}.  It measures the performance difference before/after perturbation. Smaller values are perceived to show a better robustness. 

Accuracy of the model after the SEU injection does not necessarily lower than the accuracy before the SEU injection is done, and we want to look both for positive or negative change of accuracy and penalize them equally, without cancelling each other. To do that, we will modify the formula of Accuracy Delta to be called Absolute Accuracy Delta (AAD) as follows:
    %\begin{equation} \label{eq:absaccdiff}
    %\text{Absolute Accuracy Difference} = \left|\text{Accuracy}_{after} -\text{Accuracy}_{before} \right|
    %\end{equation}
\begin{multline} \label{eq:absaccdiff}
\text{Absolute Accuracy Difference} = \\
\left|\text{Accuracy}_{after} - \text{Accuracy}_{before} \right|
\end{multline}

Softmax Difference \cite{carbone2020a} captures prediction confidence change. Robustness is $1 - \text{softmax difference}$. Previous work \cite{isrobustnessdongsu2018} highlights the trade-off between accuracy and robustness.
    \begin{multline} \label{eq:softmaxdifference}
    \text{Softmax Difference}=\\ \frac{1}{N} \sum_{j=1}^{N} 
    \Big|
    \big\langle f(\mathbf{x}_j, \mathbf{w}) \big\rangle_{p(\mathbf{w}|D)}
    -
    \big\langle f(\tilde{\mathbf{x}}_j, \mathbf{w}) \big\rangle_{p(\mathbf{w}|D)}
    \Big|_{\infty}
    \end{multline}
Our research questions rely on comparing multiple model variations and deciding which one is the best based on the given result. While both robustness measures are useful in achieving our goals on its own, interpreting both of the robustness measure at the same time may result in confusion, e.g., one model might have a low absolute accuracy difference, but a high softmax difference. To make the model comparison easier, we devise a new aggregate robustness measure called the Aggregate Robustness Index (ARIn), which takes Equation (\ref{eq:absaccdiff}) and Equation (\ref{eq:softmaxdifference}).


% \begin{equation}
% %\text{ARIn} = \sqrt{\frac{\left(\tilde{\Delta}_{\text{acc}}\right)^2 + \left(\tilde{\Delta}_{\text{softmax}}\right)^2}{2}}
% \text{ARIn} = \sqrt{\frac{\left(\text{Absolute Accuracy Difference}\right)^2 + \left({\text{Softmax Difference}}\right)^2}{2}}
% \label{eq:ARIn}
% \end{equation}

\begin{multline}
\text{ARIn} = \sqrt{\frac{\begin{aligned} &\left(\text{Absolute Accuracy Difference}\right)^2 \\ &+ \left(\text{Softmax Difference}\right)^2 \end{aligned}}{2}}
\label{eq:ARIn}
\end{multline}


%\subsection{Overview of Double Descent Experimental}

%The focus of all experiments was to understand how the test error curve manifests under different conditions when the model complexity increasing. Model complexity can be increased by three methods: increasing units in a hidden layer, adding layers, or combining both.  Each experiment with different setting varied the network width and different experiments varied the number of layers to analyse its impact on test error to explore the occurrence of the double descent phenomenon. Table \ref{tab:experiment_settings} shows the approaches taken to increase model complexity and outlines the experimental conditions for LTC, QNN, and SNN models.

%\textbf{Liquid Time-Constant Networks (LTCs)}: We utilized networks of four different depths, consisting of 1, 3, 5, and 10 layers. The width of each hidden layer was varied from 1 to 63, increasing in increments of 2.  After this, we further confirm whether the double descent phenomenon would occur by conducting additional experiments on a 5-layer LTC network but with a reduced step size of change in network width from 2 to 1. We built LTC models with multiple LTC layers followed by a fully connected output layer. Each LTC layer updates its hidden state using the following equation:
%\[
%h_t = \tau h_{t-1} + (1 - \tau) \cdot \text{ReLU}(W_{\text{in}} \cdot x_t + W_{\text{rec}} \cdot h_{t-1} + b)
%\]
%where \( h_t \) is the hidden state at time \( t \), \( \tau \) is the time constant, \( W_{\text{in}} \) and \( W_{\text{rec}} \) are input and recurrent weights, and \( b \) is the bias term. ReLU was selected for its simplicity and effectiveness.


%\textbf{Quantized Neural Networks (QNNs)}: For CIFAR-10, we varied network complexity by increase parameter c from 1 to 63 and used Adam and SGD optimizers across different training epochs (20, 50, 80), label noise (0\%, 10\%, 20\%), with and without learning rate scheduler. For MNIST, we conducted two experiments: one using the same architecture and settings as CIFAR-10 (10\% label noise, 50 training epochs, without learning rate scheduler), and a second with a simplified QNN to explore the effect of reduced network complexity. The QNN was based on a 5-layer CNN, quantised using PyTorch's Quantisation-Aware Training (QAT). 
% Two different QNN architectures were employed: one with four convolutional layers and another with two convolutional layers. For the four-layer model, the widths of the convolutional layers were set as c, 2c, 4c, 8c, while for the two-layer QNN, the widths were c, 2c. The parameter c was varied from 1 to 64 in increments of 2.

%\textbf{Spiking Neural Networks (SNNs)}: We used models with two different depths, consisting of 2 and 4 layers. The width of each hidden layer was varied from 10 to 2000, increasing in increments of 50. Our models used a two-layer and a four-layer architectures with fully connected layers followed by Leaky Integrate-and-Fire (LIF) neurons.

\section{Results}
\label{sec:results}




%\section{\uppercase{Results}}
%The summary of results of all the experiments are shown in appendix.


%\subsection{Experiment 1: Liquid Time-Constant Networks}



%In our experiments, we assessed the double descent phenomenon and generalization performance of Liquid Time-Constant (LTC) networks with varying depths (1, 3, 5, 10 layers) on the Fashion MNIST dataset. The models were trained for 50 epochs using the SGD optimizer and the \texttt{InverseSquareRootScheduler} learning rate scheduler. Training without the scheduler resulted in vanishing gradients, which prevented convergence. All LTC networks demonstrated strong generalization, with test error curves decreasing sharply as network width increased and stabilizing at consistently low values. Notably, the traditional U-shaped bias-variance curve was absent.

%Of particular note is the behavior of the 5-layer LTC network, where the test error curve exhibited a subtle indication of double descent within the hidden layer width range of 1 to 10. The training error curve for this model can be seen to reach the interpolation threshold at approximately width 5, which is about the starting point for the test error to begin its second decline. This fact is consistent with the theory of double descent which is the test error will fall a second time after the interpolation threshold reached and further increases the likelihood of a double descent occurring in the 5-layer LTC network. However, this observation remains inconclusive, as the fluctuations could be attributed to training instability rather than a definitive double descent behaviour. To investigate further, we conducted an additional experiments with narrower width increments (increasing the hidden layer width by one unit at a time). The results continued to suggest a potential double descent pattern: the test error initially decreased as the hidden layer width increased from 1 to 3, slightly increased at width 4, and then decreased again. However, due to the short duration of the error increase, it is difficult to definitively confirm the presence of double descent in the LTC network.




%Despite this ambiguity, our results imply that LTC networks may indeed be susceptible to double descent under certain conditions. However, the inherent simplicity of the Fashion MNIST dataset likely resulted in a rapid decline in training loss. This rapid convergence and the small size of the interpolation threshold may have caused the upward trend in test error to begin diminishing before it could fully manifest. Consequently, this could have obscured our ability to clearly observe the presence of the double descent phenomenon.  Future will be to evaluate double descent using the different layers and datasets.

%The ultimate goal of training a network is to minimize test error to achieve optimal generalization. While we cannot definitively claim to have observed the double descent phenomenon in LTC networks, the consistent reduction in test error across all network depths as network width increased, demonstrates that the networks ultimately achieved robust generalization performance, meeting our primary objectives.

%\begin{figure}[h]
%    \centering
%    \begin{minipage}[b]{0.40\textwidth}
%        \centering
%        (a)
%        \includegraphics[width=\textwidth]{figure2/2.png}
        % (a)\\ % Add label (a) here
%    \end{minipage}
%    \hfill
%    \\
%    \begin{minipage}[b]{0.40\textwidth}
%        \centering
%        (b)
%        \includegraphics[width=\textwidth]{figure2/15.png}
        % (b)\\ % Add label (b) here
%    \end{minipage}
%    \caption{(a) Test error curves for 1, 3, 5, and 10-layer LTC networks on Fashion MNIST using the SGD optimizer. LTCs did not show clear double descent phenomenon but generalized well. (b) Test error curves for 5-layer LTC networks on Fashion MNIST using the Adam optimizer. LTC did not show double descent phenomenon and generalized poorly.}
%    \label{fig:ltc_layers}
%\end{figure}

%In addition to using the SGD optimizer, we explored the network using the Adam optimizer. Results shown in Figure~\ref{fig:ltc_layers} (b), revealed gradient vanishing and unstable test errors, resulting in inferior generalization compared to SGD. Based on these findings, we recommend using the SGD optimizer when training LTC networks on simple datasets like Fashion MNIST. Expanding both network width and depth can improve generalization without significant risk of overfitting. Future work will examine double descent across more diverse datasets and network configurations.

%\subsection{Experiment 2: Quantised Neural Networks}


%\begin{figure}[h]
%    \centering    \includegraphics[width=0.36\textwidth]{figure2/4.png}
%    \caption{Test error curve for 4-layer QNN models with label noise (0, 0.1, 0.2) using Adam optimizer, 50 epochs, without learning rate scheduler on CIFAR-10. QNNs show significant double descent phenomenon and this trend becomes more pronounced with increasing label noise.}
%    \label{fig:qnn_1}
%\end{figure}


%\begin{figure}[h]
%    \centering
%    \includegraphics[width=0.36\textwidth]{figure2/5.png}
%    \caption{Test error curve for 4-layer QNN with 0.1 label noise using Adam optimizer, 80 epochs, without learning rate scheduler on CIFAR-10. As the number of training epochs increases, the double descent trend becomes less noticeable and the generalisation capacity decreases.
%}
%    \label{fig:qnn_2}
%\end{figure}

%In our second set of experiments, we investigated the behavior of Quantized Neural Networks (QNNs) on CIFAR-10 to understand how quantization impacts the double descent phenomenon and generalization. Initial experiments used the Adam optimizer (learning rate 0.001, 50 epochs) under varying label noise levels (0, 0.1, 0.2), as shown in Figure~\ref{fig:qnn_1}. Results confirmed that quantization does not eliminate double descent. Moreover, the data noise did not disrupt the double descent phenomenon, and increasing label noise exacerbated the double descent phenomenon.

%We then extended the training epochs to 80, keeping the learning rate at 0.001 and label noise at 0.1 (Figure~\ref{fig:qnn_2}). This weakened the magnitude of the second descent but notably worsened generalization performance, especially at higher model complexities, suggesting that network models with double descent phenomena do not always get good generalisation performance by simply increasing the model complexity and train epoch.



%To understand this decline in performance, we conducted additional experiments where we monitored both training and test errors across all epochs. The results revealed that, particularly in networks with large hidden layer widths, test error initially decreased but subsequently increased as training continued. This suggests that prolonged training (i.e., excessive training epochs) can diminish the model's generalization ability.

\iffalse
\begin{figure}[h]
    \centering
    \begin{minipage}[b]{0.4\textwidth}
        \centering
        \includegraphics[width=\textwidth]{figure2/6.png}
        % (a)\\ % Add label (a) here
    \end{minipage}
    \hfill
    \begin{minipage}[b]{0.4\textwidth}
        \centering
        \includegraphics[width=\textwidth]{figure2/7.png}
        % (b)\\ % Add label (b) here
    \end{minipage}
    \caption{(a) Test error curve for 4-layer QNN with 0.1 label noise using Adam optimizer with higher learning rate and 20 epochs, without learning rate scheduler on CIFAR-10. The test error curve
    did not show double descent trend and was highly unstable, indicating weak generalization ability. (b) Test error curve for 4-layer QNN with 0.1 label noise using Adam optimizer with higher learning rate and 20 epochs, with learning rate scheduler on CIFAR-10. The test error curve did not show double descent trend but demonstrated strong generalization ability.}
    \label{qnn_4}
\end{figure}
\fi



%To address this, we reduced training epochs to 20 and increased the learning rate to 0.1. However, the larger learning rate introduced instability due to exploding gradients, as shown in Figure~\ref{qnn_4}(a). Incorporating the StepLR scheduler (Figure~\ref{qnn_4}(b)) effectively eliminated the double descent phenomenon and improved generalization performance.


%To explore whether the choice of optimizer affects the double descent phenomenon, we replaced Adam with SGD under the same conditions. Without a scheduler, the double descent pattern persisted but remained unstable  (Figure~\ref{qnn_5}(a)). Adding the StepLR scheduler stabilized the test error curve and improved generalization (Figure~\ref{qnn_5}(b)). Overall, SGD outperformed Adam in generalization performance but did not entirely eliminate double descent.

%\begin{figure}[h!]
%    \centering
%    \includegraphics[width=0.4\textwidth]{figure2/10.png}
%    \caption{Test error curve for 2- and 4-layer QNNs with 0.1 label noise using Adam optimizer, 50 epochs, without learning rate scheduler on MNIST. The 2-layer model exhibited the double descent phenomenon, whereas the 4-layer model did not.}
%    \label{fig:qnn_3}
%\end{figure}

%In the two experiments conducted on the MNIST dataset, the results are illustrated in Figure~\ref{fig:qnn_3}. The blue line represents the performance of the original QNN model which does not show the double descent behaviour, while the red line depicts the results from the simplified QNN model which show the double descent behaviour. When comparing the blue line with the corresponding results in Figure~\ref{fig:qnn_3}, it becomes evident that when the experimental model and configuration remain unchanged, the use of a simpler dataset allows the model to learn the underlying patterns more effectively, leading to better generalization performance. However, although the simplified QNN model exhibited the double descent phenomenon, as shown by the red line, the generalization performance of the model was worse than complex model.



%Through these experiments, we have established several important conclusions and recommendations for training QNNs:

%1. \textbf{Persistence of Double Descent}: quantisation does not negate the occurrence of the double descent phenomenon. The characteristic double descent behavior observed in the original CNN architecture persists in the quantised version, confirming that the quantisation process alone does not eliminate this phenomenon.

%2. \textbf{Data Noise}:Increased data noise may exacerbate the double descent phenomenon

%3. \textbf{Mitigating Overfitting}: Several techniques, including reducing the learning rate, introducing a learning rate scheduler, and limiting the number of training epochs, effectively mitigate overfitting. In particular, the learning rate scheduler helps to eliminate the double descent phenomenon without affecting the generalisation ability of the model.

%4. \textbf{Model Complexity vs. Generalization}: Our experiments indicate that simply increasing model complexity does not guarantee improved generalization. Even if there is a double descent behaviour, if the learning rate and training duration are not adjusted properly, the test error during the second descent may not drop below the U-curve lowest point or, in some cases, the traditional U-shaped curve may replace the double descent curve entirely. In such instances, the generalization error continues to rise after reaching its initial minimum, with no subsequent decline.

%5. \textbf{Optimizer and scheduler}: The choice of optimizer and the use of a learning rate scheduler are crucial factors influencing the behavior of QNNs. In our experiment, SGD is better than Adam Optimizer but both can lead to unstable training outcomes and a noticeable decline in generalization performance, particularly under conditions of high learning rates and extended training durations. However, when combined with a learning rate scheduler, the model can provide a more stable and robust generalization performance.



%In conclusion, when training QNNs, particularly on complex datasets like CIFAR-10, we must pay careful attention to the selection of optimizer, learning rate, training duration, and the use of learning rate schedulers. These factors play a crucial role in influencing the model's generalization performance and the occurrence of the double descent phenomenon. Future research should continue to explore these variables in order to further optimize the training of more type of QNNs.

\subsection{Training Result}

\begin{table}[h]
\centering
\begin{tabular}{c|rrrrrrr|r}
\toprule
Model Variant  & DNN & BNN \\
\midrule
Model 0 &   \B{0}{.9757} &  \B{0}{.8288} \\
Model 1 &   \B{0}{.9752} &  \B{1}{.8332} \\
Model 2 &   \B{1}{.9763} &  \B{0}{.8276} \\
Model 3 &   \B{0}{.9748} &  \B{0}{.8243} \\
\midrule
Average &   \B{0}{.9755} &  \B{0}{.8285} \\
\bottomrule
\end{tabular}
\caption{Test accuracy of Bayesian and Deterministic CNN Models grouped by model variants}
\label{table:activation_performance_deterministic_test}
\end{table}

Based on our findings in this chapter, \textbf{Bayesian CNN model is showing an inferior performance in test accuracy} compared to the deterministic CNN model. This observation is consistent across all activation functions, model variants, and Bayesian model activation function. The average accuracy of test data that can be achieved by the deterministic CNN model is 97.55\%, while the Bayesian CNN model can only achieve 82.85\%. it's suspected that the performance difference is inherent to whether we are implementing Bayesian or deterministic CNN model. Here are some factors that can possibly result in this observation:
\begin{enumerate}%[noitemsep, topsep=0pt]
%\item All weights and bias in the Bayesian CNN neural network are initialized with a similar prior parameter ($a$ and $b$), assuming that all weight and bias' distributions are following the same distribution. 
\item All weights and bias in the Bayesian CNN neural network are initialized with the same prior distribution parameter ($a$ and $b$).
\item Optimization of Bayesian Neural Networks with ELBO loss is harder compared to cross-entropy in deterministic CNN model. KL divergence in ELBO loss can pull the variational posterior distribution away from the optimal point to stay closer to the prior (posterior collapse).
\end{enumerate}

In the subsequent subchapters that measure robustness of deterministic CNN models and Bayesian CNN models, most table will provide the aggregate value of SEU robustness measure by grouping the experiment result based on several criteria and then averaging them. For simplicity, we denote these aggregated metrics as Average Absolute Accuracy Difference (AAAD) and Average Softmax Difference (ASD). Aggregated Robustness Index (ARIn) will be calculated based on these two values. The lowest value for each column will be typed in bold.

%(e.g. model variation, activation function)

\subsection{Robustness by model type (BNN vs DNN)}

\iffalse
\begin{figure}[H]
    \centering
    %\begin{minipage}[b]{0.4\textwidth}
    %    \centering
    %    \includegraphics[width=\textwidth]{figure2/6.png}
    %    % (a)\\ % Add label (a) here
    %\end{minipage}
    %\hfill
    \begin{minipage}[b]{0.4\textwidth}
        \centering
        \includegraphics[width=\textwidth]{robboth2.png}
        % (b)\\ % Add label (b) here
    \end{minipage}
    \caption{Robustness of Deterministic and Bayesian CNN model}
    \label{fig:robboth2}
\end{figure}
\fi

\begin{figure*}[t]
    \centering
    \includegraphics[width=0.7\textwidth]{robboth2.png}
    \caption{Robustness of Deterministic and Bayesian CNN model}
    \label{fig:robboth2}
\end{figure*}

%Summary of absolute accuracy difference and softmax difference Deterministic vs Bayesian CNN
\begin{table}[h]
\centering
\begin{tabular}{lrrr}
\toprule
Model Type & AAAD & ASD & ARIn \\
\midrule
DNN & \B{1}{.0149} & \B{0}{.2696} & \B{0}{.1909} \\
BNN & \B{0}{.0244} & \B{1}{.1074} & \B{1}{.0779} \\
\bottomrule
\end{tabular}
\caption{Summary of average absolute accuracy difference and average softmax difference Deterministic vs Bayesian CNN}
\label{table:robbothcompare}
\end{table}

\subsubsection{Robustness of Deterministic CNN models}

From Figure \ref{fig:robboth2}, we can infer that the impact of SEU on Absolute Accuracy Difference and Softmax Difference exhibits an approximately linear relationship, showing a clear correlation between the two performance metrics. This observation is similar from what has been observed from \cite{carbone2020a}, in which the author spotted a trade-off between pure accuracy and robustness (1 - softmax difference) for high performing deterministic networks against adversarial attack (FGSM).

\subsubsection{Robustness of Bayesian CNN models}

Contrary to the Deterministic CNN model SEU injection result that has a clear correlation between the two performance metric in a linear straight line, the performance shown by the Bayesian CNN model does not exhibit a same pattern. In Bayesian CNN Model SEU Injection from the figure above, some noticeable points make up a linear correlation of absolute accuracy difference and softmax difference, but there are also observations that does not support this relationship. Some SEU Injection results show high softmax difference, while still keeping the absolute accuracy difference low. 

\subsubsection{Robustness of Deterministic CNN models vs Bayesian CNN Models}

Based on Table \ref{table:robbothcompare}, it is observed that the deterministic CNN model is showing more robustness against SEU injection in absolute accuracy difference. It shows that the average absolute accuracy difference if a deterministic CNN model is attacked by SEU injection is 1.49\%, while the Bayesian CNN model is expected to have an absolute accuracy difference of 2.44\%. 

This phenomenon can be triggered by the way that Bayesian Neural Network is sampling the weight/bias distribution parameter from the guide. In deterministic CNN model, the SEU only change one fixed weight/bias. In Bayesian CNN Model, as the distribution parameter is perturbed by the SEU, all Monte Carlo (MC) samples are drawn from the same perturbed distribution parameter. This makes every MC sample to be perturbed in the same way. This, will in turn disturb the average logits that is needed to classify the data with softmax. The resulting condition would be the data that is more prone to misclassification due to disturbed sampled weight and bias.

On the other hand, in terms of average softmax difference, Bayesian CNN model can be seen to demonstrate a better robustness against SEU injection. Softmax difference, which indicates the change in the predicted class probability, is observed to be 0.2696 on the Deterministic CNN Model, and 0.1074 on the Bayesian CNN Model. This means that Bayesian CNN Model has a smaller shift in output probabilities compared to the Deterministic one. 

The observed condition can happen due to how logits are calculated after MC sampling. As we run over the network 10 times, or as much as the number of samples is determined, logit shifts from the first run might be diluted by the subsequent run, as we are taking the average of all 10 runs to be the input of the softmax, that will classify our data. Although the chance of misclassifying some data sample is increased, Bayesian CNN model offer a more stable prediction probability against SEU.

To conclude our comparison, we employ the Aggregate Robustness Measure (ARIn) defined in Equation (\ref{eq:ARIn}). Based on this metric, we find that \textbf{Bayesian Neural Network exhibits greater robustness than the Deterministic Neural Network}, as evidenced by its smaller ARIn value (0.0779 compared to 0.1909 for the latter). In the applications where class confidence or predictive probability constitutes a critical factor, the Bayesian CNN model is therefore more suitable to use compared to the deterministic CNN model. Some applications can include medical imaging (e.g. If the model uncertainty exceeds the threshold, the case is "rejected" or flagged for human evaluation \cite{10.3389/fninf.2024.1346723} ) and ensemble modeling (in which multiple models are combined to produce aggregate predictions \cite{bishopbarber98}). 


\subsection{Robustness by model variation}

%Summary of absolute accuracy difference and softmax difference by model variant for Deterministic and Bayesian CNN
\begin{table}[h]
\centering
\begin{tabular}{l|rr|rr|rr}
\toprule
 & \multicolumn{2}{c|}{AAAD} & \multicolumn{2}{c|}{ASD} & \multicolumn{2}{c}{ARIn} \\
 & DNN & BNN & DNN & BNN & DNN & BNN \\
\midrule
0 & \B{0}{.0145} & \B{0}{.0243} & \B{0}{.2694} & \B{0}{.1106} & \B{0}{.1908} & \B{0}{.0800} \\
1 & \B{1}{.0142} & \B{0}{.0254} & \B{0}{.2699} & \B{0}{.1109} & \B{0}{.1911} & \B{0}{.0805} \\
2 & \B{0}{.0157} & \B{1}{.0221} & \B{0}{.2702} & \B{1}{.0944} & \B{0}{.1914} & \B{1}{.0685} \\
3 & \B{0}{.0152} & \B{0}{.0257} & \B{1}{.2687} & \B{0}{.1138} & \B{1}{.1903} & \B{0}{.0825} \\
\bottomrule
\end{tabular}
\caption{Robustness of Deterministic and Bayesian CNN model by model variant}
\label{table:robbothmodelgran}
\end{table}

In Table \ref{table:robbothmodelgran}, each row represents model variations explained in Subsection \ref{subs:modelvar}.

\subsubsection{Robustness of Deterministic CNN models}

Table \ref{table:robbothmodelgran} summarizes these results by averaging the Absolute Accuracy Difference (AAAD) and Softmax Difference (ASD) across model variants. Based on the Aggregate Robustness Index (ARIn), \textbf{Model 3 (trained with weight decay) demonstrates the strongest robustness}. Although it exhibits slightly worse AAAD compared to the base model (Model 0), it achieves a noticeably lower ASD. This is consistent with the expectation that weight decay reduces the magnitude of the weights, leading to a smaller logits and hence smaller softmax differences.

In comparison to prior work \cite{dennispoperobustcnn}, our findings diverge. While it suggests that a variant equivalent to our Model 2 achieved the best AAAD, our experiments indicate that Model 2 has the least robust AAAD, while Model 1 (with smartpooling) performs best in AAAD.

\subsubsection{Robustness of Bayesian CNN models}

As shown in Table \ref{table:robbothmodelgran}, \textbf{Model 2 (trained with dropout) is the most robust among the Bayesian variants}, exhibiting the lowest values for both AAAD and ASD. Although AAAD is generally similar across Bayesian variants, the reduced ASD of Model 2 makes its particularly notable. This outcome aligns with findings in \cite{dennispoperobustcnn}, despite differences in experimental setup.

\subsubsection{On the robustness of Bayesian CNN model trained with dropout layer}

The significant robustness of Model 2 compared to other model variants in the Bayesian setting can be attributed by the nature of Dropout layers that demonstrate redundancy during training. As some neurons are inactive during training, our neural network can not just too reliant on any single neuron to carry critical information/infer a feature. In our Bayesian CNN model without dropout layer, there might be a chance that a SEU/bitflips can hit a dominant neuron and shift our prediction significantly. With a regularized model with dropout, the resulting model has no dominant neuron, so during inference, other neurons are as contributing as others, resulting in less shift in prediction if one of the neurons is being hit by bitflip.

\subsection{Robustness prior and prior parameter}

%Summary of absolute accuracy difference and softmax difference by prior and prior parameter for Bayesian CNN
\begin{table}[h]
\centering
\begin{tabular}{l|rr|rr|rr}
\toprule
 & \multicolumn{2}{c|}{AAAD} & \multicolumn{2}{c|}{ASD} & \multicolumn{2}{c}{ARIn} \\
 & C & S & C & S & C & S \\
\midrule
G & \B{1}{.0177} & \B{0}{.0293} & \B{1}{.0794} & \B{0}{.1281} & \B{1}{.0575} & \B{0}{.0929} \\
L & \B{0}{.0193} & \B{1}{.0287} & \B{0}{.0823} & \B{0}{.1306} & \B{0}{.0598} & \B{0}{.0946} \\
U & \B{0}{.0231} & \B{0}{.0303} & \B{0}{.1054} & \B{1}{.1277} & \B{0}{.0763} & \B{1}{.0928} \\
\bottomrule
\end{tabular}
\caption{Robustness of Bayesian CNN model by prior and prior parameter}
\label{table:robbaypriorpriorparam2}
\end{table}

In Table \ref{table:robbaypriorpriorparam2}, each row represents each prior distribution; which are Gaussian (G), Laplace (L), and Uniform (U). For each performance measure, we varied the SEU injection by the Center (C) and Scale (S) prior parameter.

\subsubsection{Which prior is the most robust?}

For the absolute accuracy difference, both Gaussian and Laplace prior have a similar result on the center and spread, with the Uniform distribution is the worst among all prior distributions. In terms of softmax difference, the Gaussian and Laplace distribution also performs similarly on the prior center and are both better compared to the Uniform distribution. A different observation can be found on the spread, which shows the Uniform and Gaussian perform similarly and better compared to the Laplace distribution. Overall, we can conclude that the \textbf{Gaussian prior is the most robust prior} among all, with Laplace prior shown a better accuracy robustness on prior spread attack, and the uniform distribution shown a better softmax robustness on prior spread attack.

\subsubsection{Gaussian prior vs Laplace prior}

%TODO explain why in general Gaussian is better

%Each prior distribution have a different Probability Distribution Function (PDFs). The Gaussian distribution has lighter tails compared to the Laplace distribution with the same parameter. As we shift the mean for both distribution, we also pull the tail samples toward or accross the decision boundary. Because Laplace distribution have a heavier tail, the probability of tail samples to get sampled is higher, thus increasing the likelihood of prediction changes. The Uniform distribution is unsurprisingly doing worst, as the probability density of Uniform distribution is flat, making the samples evenly spread across the specified range. A shift in $a$ value for uniform distribution move a larger fraction of samples across the decision boundary, with the same probability of being sampled, making the prediction accuracy worse.

%Laplace distribution is less affected in absolute accuracy difference if the spread parameter is being attacked compared to the Gaussian distribution, due to the fact that Gaussian distribution is more concentrated in the center. If we increase or do bitflip on the scale parameter for Gaussian distribution, the distribution of sample near the center widen and possibly move the sampled weight/bias across the decision boundary. For the Laplace distribution, modifying the scale parameter does not change the sampled value much, as the probability density of Laplace is already wide relative to the Gaussian prior. 

Gaussian priors have lighter tails, making them more sensitive to mean or scale shifts that move samples across decision boundaries and increase prediction changes. Laplace priors, with heavier tails, are less affected by spread perturbations since their distributions are already wide, Uniform priors perform worse because their flat density causes any shift to affect many samples simultaneously, degrading accuracy.

\subsubsection{Uniform prior robustness in softmax difference against the spread SEU injection}

Since changing the scale parameter for the Gaussian and Laplace prior affect the peak height and decay rate in the tails, if we do Monte Carlo sampling for weight/bias sampling, the difference of sample values is expected to be higher compared to the Uniform prior that have a flat PDF. This is due to the fact that the high-density center is now less dense and mass shifts into the tails, which will amplify the variability in logits. The uniform distribution can yield a better softmax difference, due to the fact that the each value of weight/bias have a similar probability to be sampled. A slight change in Uniform's scale parameter, for a latent weight $x$, will result in a smaller change of probability, compared to the change for the same latent weight $x$ in Gaussian/Laplace prior.

\subsubsection{On the robustness of prior center vs prior spread against SEU}

%TODO add explanation. Why the center parameter is more robust than the spread

Based on the Table \ref{table:robbaypriorpriorparam2} above, we can also retrieve a conclusion about how the SEU interacts with the center prior parameter and spread prior parameter. For all metrics, \textbf{the center prior parameter is consistently more robust than the spread prior parameter}. Robustness of Bayesian Neural Networks against SEU injection on the prior center parameter comes from the nature of shifting the mean of the prior that moves all the weight samples by the same absolute amount, preserving the relative positions and allowing them to cancel out during the Monte Carlo averaging. Shift in the spread parameter will change the weight sampled to be significantly different from the mean (either pushing some further into the tails or other closer to the center), making the resulting logits to be much different than before the SEU injection happened.

\subsection{Robustness by activation function}

%bayesian vs deterministic per metric
\begin{table}[h]
\centering
\begin{tabular}{l|rr|rr|rr}
\toprule
& \multicolumn{2}{c|}{AAAD} & \multicolumn{2}{c|}{ASD} & \multicolumn{2}{c}{ARIn} \\
 & DNN & BNN & DNN & BNN & DNN & BNN \\
\midrule
R & \B{1}{.0078} & \B{1}{.0168} & \B{0}{.2651} & \B{0}{.0958} & \B{0}{.1875} & \B{0}{.0688} \\
W & \B{0}{.0178} & \B{0}{.0213} & \B{0}{.2716} & \B{0}{.0964} & \B{0}{.1924} & \B{0}{.0698} \\
ru & \B{0}{.0234} & \B{0}{.0379} & \B{0}{.2835} & \B{0}{.1626} & \B{0}{.2012} & \B{0}{.1180} \\
r6 & \B{0}{.0114} & \B{0}{.0248} & \B{0}{.2653} & \B{0}{.1012} & \B{0}{.1878} & \B{0}{.0737} \\
sg & \B{0}{.0136} & \B{0}{.0190} & \B{1}{.2598} & \B{1}{.0800} & \B{1}{.1839} & \B{1}{.0581} \\
sn & \B{0}{.0141} & \B{0}{.0294} & \B{0}{.2702} & \B{0}{.1056} & \B{0}{.1913} & \B{0}{.0775} \\
th & \B{0}{.0163} & \B{0}{.0214} & \B{0}{.2714} & \B{0}{.1102} & \B{0}{.1923} & \B{0}{.0794} \\
\bottomrule
\end{tabular}
\caption{Robustness of Deterministic and Bayesian CNN model by activation function}
\label{table:robbothact}
\end{table}

For Table \ref{table:robbothact}, we code each activation function as follows: RWG (R), WG (W), relu (ru), relu6 (r6), sigmoid (sg), sin (sn), tanh (th). We can observe that the Rectified Weighted Gaussian activation function has the most robust absolute accuracy difference among all activation function for the deterministic CNN model (which has been proven in \cite{dennispoperobustcnn}) and Bayesian CNN Model. On our newly introduced softmax difference measure, the sigmoid activation function achieves the lowest Average Softmax Difference (ASD) in both settings.

%TODO why the RWG might be the best

Relu, which is the most expressive activation function that we have in this experiment (as seen from Figure \ref{fig:activationfunctions}), shown a comparably better test accuracy despite showing the worst robustness. RWG, which yields the lowest test accuracy, achieves the strongest robustness, especially in terms of absolute accuracy difference. This observation shown us a symptom of a \textbf{trade-off between test accuracy and SEU robustness}, which is also observed in the work of \cite{carbone2020a}.

It can also be seen from Table \ref{table:robbothact} that the sigmoid activation function have the second best robustness in absolute accuracy difference on BNNs, while also having the best softmax difference. SEU happens to weight and bias before activation function is implemented. Let's get back on how the neuron is calculated in the feed-forward process. This process, as mentioned in \cite{bishoppatternrecognition2007} is notated as:

\begin{equation}
\label{eq:neuronfeedforward}
a_j = \sum_{i=1}^{D} w^{(1)}_{ji} x_i + w^{(1)}_{j0}
\end{equation}

%page 227
where $x_1,x_2,\dots,x_D$ are inputs, $j=1,2,\dots,M$ and the superscript (1) represents the parameters in the first 'layer' of the network, which have a 'width' of $M$. $w^{(1)}_{j0}$ is  Each of $a_j$ is then transformed in the following equation:
\begin{equation}
z_j = h(a_j)
\end{equation}
where $h(\cdot)$ is a differentiable activation function.
In order to understand how different activation function can yield different robustness, let's compare the sigmoid and tanh activation functions. Sigmoid has a smaller gradient $h'()$ near zero (0.25) than tanh (1), meaning that weight perturbations around this region result in a smaller output change. This explains the sigmoid's higher robustness under SEU injection.



%$w^{(1)}_{ji}$ which represents weight, or $w^{(1)}_{j0}$ which represents bias, are the component in this feed-forward process that is prone to the SEU injection. As one of both values changed into SEU injected value $a'_j$, we would like to know how this change can effect the value of $z$ into $z'_j$, the lower the change, the difference in softmax value will be also lower.

%In this experiment, the weight/bias is expected to be near zero with our setting of prior parameter $a=0$ and $b=1$. The majority of bitflips that we do will result in a change of weight/bias that is near zero, when we hit bit index that is not 0 (sign bit) and 1 (first exponent bit). Therefore, let's see the impact of changing $a$ to $z$ by observing the differential of our activation function $h'()$. We will take a look into sigmoid and the seemingly similarly shaped tanh, and its differential near 0.

%$$
%\sigma(z) = \frac{1}{1 + e^{-z}} ,\qquad\frac{d\sigma(z)}{dz} = \sigma(z)\big(1 - \sigma(z)\big) ,\qquad
%$$
%$$
%\sigma(0) = 0.5 \quad \Rightarrow \quad %\sigma'(0) = 0.25
%$$

%$$
%\tanh(z) = \frac{e^z - e^{-z}}{e^z + e^{-z}} ,\qquad \frac{d}{dz}\tanh(z) = 1 - \tanh^2(z) ,\qquad
%$$

%$$
%\tanh(0) = 0 \quad \Rightarrow \quad 5\tanh'(0) = 1
%$$

%From our calculation of diferential for both activation functions above, we can see that the differential on sigmoid at 0 is 0.25, while the tanh is at 1. A slight change in input for tanh have a bigger shift in the value after activation function. Therefore, we can conclude that the sigmoid function gained robustness in SEU because of a smaller gradient near $z=0$. 

%It will also be interesting to look at the gradient of our activation function on extreme values, especially in looking in the effect of SEU injection on bit index 1 (first exponent). \cite{bishoppatternrecognition2007} mentioned that the sigmoid decay asymptotically like $e^{-z}$ for $z \to +\infty$, while tanh, which has $tanh'(z) = 1- tanh^2(z)$, will decay in $4e^{-2z}$ for $z \to +\infty$, which is faster than the sigmoid. We can assume that the tanh function will dampen the $z$ value better on the extreme values.

\subsection{Robustness by layer and layer parameter}

%Summary of absolute accuracy difference and softmax difference by bit layer and module for Deterministic and Bayesian CNN
\iffalse
\begin{table}[H]
\centering
\begin{tabular}{ll|rr|rr|rr}
\toprule
 &  & \multicolumn{2}{c|}{AAAD} & \multicolumn{2}{c|}{ASD} & \multicolumn{2}{c}{ARIn} \\
 &  & DNN & BNN & DNN & BNN & DNN & BNN \\
\midrule
conv1 & bias & \B{0}{.0048} & \B{0}{.0095} & \B{0}{.2651} & \B{1}{.0641} & \B{0}{.1875} & \B{1}{.0458} \\
conv1 & weight & \B{0}{.0070} & \B{0}{.0129} & \B{0}{.2696} & \B{0}{.0762} & \B{0}{.1907} & \B{0}{.0546} \\
\midrule
conv2 & bias & \B{1}{.0021} & \B{1}{.0089} & \B{1}{.2622} & \B{0}{.0723} & \B{1}{.1854} & \B{0}{.0515} \\
conv2 & weight & \B{0}{.0058} & \B{0}{.0136} & \B{0}{.2640} & \B{0}{.0987} & \B{0}{.1868} & \B{0}{.0704} \\
\midrule
fc1 & bias & \B{0}{.0326} & \B{0}{.0459} & \B{0}{.2773} & \B{0}{.1720} & \B{0}{.1974} & \B{0}{.1259} \\
fc1 & weight & \B{0}{.0372} & \B{0}{.0555} & \B{0}{.2792} & \B{0}{.1611} & \B{0}{.1991} & \B{0}{.1205} \\
\bottomrule
\end{tabular}
\caption{Robustness of Deterministic and Bayesian CNN model by layer and layer parameter}
\label{table:robbothlayermodule}
\end{table}
\fi


\begin{table}[h]
\centering
\begin{tabular}{ll|rr|rr}
\toprule
 &  & \multicolumn{2}{c|}{AAAD} & \multicolumn{2}{c}{ASD} \\
 &  & DNN & BNN & DNN & BNN \\
\midrule
conv1 & bias & \B{0}{.0048} & \B{0}{.0095} & \B{0}{.2651} & \B{1}{.0641}  \\
conv1 & weight & \B{0}{.0070} & \B{0}{.0129} & \B{0}{.2696} & \B{0}{.0762} \\
\midrule
conv2 & bias & \B{1}{.0021} & \B{1}{.0089} & \B{1}{.2622} & \B{0}{.0723} \\
conv2 & weight & \B{0}{.0058} & \B{0}{.0136} & \B{0}{.2640} & \B{0}{.0987} \\
\midrule
fc1 & bias & \B{0}{.0326} & \B{0}{.0459} & \B{0}{.2773} & \B{0}{.1720} \\
fc1 & weight & \B{0}{.0372} & \B{0}{.0555} & \B{0}{.2792} & \B{0}{.1611} \\
\bottomrule
\end{tabular}
\caption{Robustness of Deterministic and Bayesian CNN model by layer and layer parameter (AAAD \& ASD)}
\label{table:robbothlayermodule}
\end{table}

\begin{table}[h]
\centering
\begin{tabular}{ll|rr}
\toprule
 & & \multicolumn{2}{c}{ARIn} \\
 &  & DNN & BNN \\
\midrule
conv1 & bias & \B{0}{.1875} & \B{1}{.0458} \\
conv1 & weight &  \B{0}{.1907} & \B{0}{.0546} \\
\midrule
conv2 & bias & \B{1}{.1854} & \B{0}{.0515} \\
conv2 & weight & \B{0}{.1868} & \B{0}{.0704} \\
\midrule
fc1 & bias & \B{0}{.1974} & \B{0}{.1259} \\
fc1 & weight & \B{0}{.1991} & \B{0}{.1205} \\
\bottomrule
\end{tabular}
\caption{Robustness of Deterministic and Bayesian CNN model by layer and layer parameter (ARIn)}
\label{table:robbothlayermodule}
\end{table}
    %\end{minipage}
    %\caption{Robustness of Deterministic and Bayesian CNN model by attacked bit location}
    %\label{fig:robbothbitlocbig}


Table \ref{table:robbothlayermodule} summarizes the SEU Injection result by averaging the value of Absolute Accuracy Difference and Softmax Difference across layers and layer parameters of the CNN architecture. Based on that table, it is observed that convolutional layer 2 is the most robust, followed by convolutional layer 1. The fully connected layer is the least robust, as it exhibits the highest mean of softmax difference and the highest mean absolute accuracy difference. These findings also align with the result shown in \cite{dennispoperobustcnn}, where the further convolutional layer (conv2 and conv3) have the lowest average test delta and the fully connected layer has the highest.

\subsubsection {Layer SEU Robustness, propagation depth}

%TODO explain how error propagates through the layers. Error in early layers are likely to be 'dampened' by the activation function etc. Error on the far end of the layer will not be dampened well and captured by the prediction layer.

%Fully connected layer, which is the furthest layer we have in our architecture, is the least robust due to how the SEU injection impacts the final logits. If we do the SEU injection early in the convolutional 1 layers, there are multiple occassion that nonlinear process will dampen the effect of SEU. 

%First, every convolution layer will be followed by an activation function. Nonlinear activation function can help reduce the impact of SEU or bitflips. After the activation function, the result is taken into pooling layer, in which the multiple values can are aggregate into a smaller dimension. If the pooling layer is utilizing the average pooling layer, the value of the SEU will be averaged with other values, reducing the impact. In our case study, we are using max pooling. The value of the SEU might just be propagated to the next layer if the resulting value is the highest among all. Therefore, as the work in \cite{dennispoperobustcnn} suggested, smart pooling is used, to remove values that are above specified threshold, especially if the value was hit by SEU.

%Let's compare this condition with fully connected layer from our model architecture. In this layer, the SEU injected values will not undergo any non-linear transformation like what we had in the convolutional layers. Therefore, this value will directly impact the calculation of logits and class probability, thus explaining the higher absolute accuracy difference and softmax difference.


%%The illustration on how SEU injection effects can propagate through the network and its impact is shown in Figure \ref{fig:errorpropillust}. An SEU injected in early convolutional layers, as shown in Figure \ref{fig:errorpropillust} (a) is symbolized as red lightning to show its devastating impact. The subsequent impact is marked yellow to show that SEU effect is weakened, followed by small green lightning on the FC1 layer, showing that the impact of SEU is further diminished. The sinusoidal wave in the neuron is intended to show that the attack undergo non-linear transformation, dampening the impact of the SEU injected value. The colored cloud at the end of the network represents the impact of SEU on the logits, with red means that the logits is impacted the worst (as shown in Figure \ref{fig:errorpropillust} (b)), and green is the least.

%Figure \ref{fig:errorpropillust} illustrates how the effects of an SEU injection propagate through a CNN model. In panel (a), an SEU injected into an early convolutional layer (indicated by the red lightning bolt) undergoes attenuation as it passes through subsequent layers. This weakening is shown by yellow and then smaller green lightning symbols in later layers, reflecting the diminishing influence of the fault. The sinusoidal symbols within the neurons represent non-linear transformations, which further dampen the perturbation. Consequently, the effect on the final logits is minimal, represented by the green cloud at the output.

%In contrast, panel (b) shows an SEU injected directly into the fully connected (FC1) layer. Here, the perturbation is not mitigated by earlier non-linear transformations and instead propagates strongly to the output layer, producing the most severe distortion in the logits, indicated by the red clouds.

The fully connected layer is the least robust because SEU injections directly affect the final logits without any subsequent non-linear transformations, resulting in the highest AAAD and ASD. In contrast, faults injected in earlier convolutional layers are partially mitigated by activation and pooling operations, which dampen or average out perturbations before they reach the output. As illustrated in Figure \ref{fig:errorpropillust}, early-layer faults weaken as they propagate through the network, whereas faults at the fully connected layer propagate directly to the output, causing the most severe distortions.

\subsubsection{Layer component SEU Robustness, weight vs bias}

Based on the result of our experiment, we can see that the weight is less robust compared to the bias, when we inject them with SEU/biftlips. This is due to the nature of the neural networks itself. Looking back at Equation (\ref{eq:neuronfeedforward}), we can see that the weight is multiplied by the input value $x_i$ to compute $a_j$. \textbf{Change in weight to the value of $a_j$ is amplified by the $x_i$}, compared to the change in bias that is merely a constant.



\begin{figure}[h]
\caption{Illustration of SEU injection and its propagated effect on conv1 layer (a) and fc1 layer (b) in a CNN model}
\includegraphics[scale=0.45]{images/error-propagation-illustration-stacked.png}
\centering
\label{fig:errorpropillust}
\end{figure}

\subsection{Robustness by attacked bit location} \label{robattackedbit}

%Summary of absolute accuracy difference and softmax difference by bit location for Deterministic and Bayesian CNN


\begin{table}[h]
\centering
\begin{tabular}{l|rr|rr}
\toprule
 & \multicolumn{2}{c|}{ALAD} & \multicolumn{2}{c}{ARIn}\\
 & DNN & BNN & DNN & BNN\\
\midrule
0 & \B{0}{-1.1804} & \B{0}{-0.7909}& \B{0}{.1848} & \B{0}{.0381}\\
\midrule
1 & \B{0}{37.0504} & \B{0}{37.8829}&\B{0}{.2371} & \B{0}{.3009}\\
3 & \B{0}{-1.4815} & \B{0}{-0.6454}& \B{1}{.1847} & \B{0}{.0376}\\
6 & \B{0}{-0.5849} & \B{0}{-0.4291} & \B{0}{.1850} & \B{0}{.0378}\\
\midrule
10 & \B{0}{-2.2308} & \B{0}{-1.3945}& \B{1}{.1847} & \B{1}{.0374}\\
15 & \B{0}{-3.7360} & \B{0}{-2.8997}& \B{1}{.1847} & \B{1}{.0374}\\
21 & \B{1}{-5.5422} & \B{1}{-4.7058}& \B{1}{.1847} & \B{0}{.0376}\\
\bottomrule
\end{tabular}
\caption{Robustness of Deterministic and Bayesian CNN model by attacked bit location (ALAD \& ARIn)}
\label{table:robbothbitlocaladonly}
\end{table}

\begin{table}[h]
\centering
\begin{tabular}{l|rr|rr}
\toprule
  & \multicolumn{2}{c|}{AAAD} & \multicolumn{2}{c}{ASD}  \\
  & DNN & BNN & DNN & BNN \\
\midrule
0 &  \B{0}{.0005} & \B{0}{.0075} & \B{0}{.2613} & \B{0}{.0533} \\
\midrule
1 & \B{0}{.1023} & \B{0}{.1206} & \B{0}{.3193} & \B{0}{.4080} \\
3 & \B{0}{.0002} & \B{0}{.0070} & \B{0}{.2612} & \B{0}{.0527}  \\
6 & \B{0}{.0013} & \B{0}{.0070} & \B{0}{.2617} & \B{0}{.0529} \\
\midrule
10 & \B{0}{.0001} & \B{1}{.0067} & \B{1}{.2611} & \B{1}{.0525} \\
15 & \B{1}{.0000} & \B{1}{.0067} & \B{1}{.2611} & \B{1}{.0525} \\
21 & \B{1}{.0000} & \B{1}{.0067} & \B{1}{.2611} & \B{0}{.0528} \\
\bottomrule
\end{tabular}
\caption{Robustness of Deterministic and Bayesian CNN model by attacked bit location (AAAD \& ASD)}
\label{table:robbothbitloc}
\end{table}

To analyze the effect of SEU injections at different bit positions, we compute the log absolute difference of values before and after an SEU, defined as:
\begin{multline}
\text{Log Absolute Difference} = \\ \log_{10}(|\text{value after SEU}-\text{value before SEU}|)
\end{multline}

%Attacked bit index in this experiment are sign (index 0), exponent (index 1,3,6), and mantissa (index 10, 15, 21). Based on Table \ref{table:robbothbitloc}, the average log absolute difference of attack against bit index 1 is around 37, which is is the worst among all SEU injection bit index. We can see that bit index 1 is the significantly least robust in terms of AAAD, ASD, which consequently shows the worst ARIn. 

In this experiment, the attacked bit index in this experiment are sign (index 0), exponent (index 1,3,6), and mantissa (index 10, 15, 21). As shown in Table \ref{table:robbothbitloc}, flipping bit index 1 yields an average log absolute difference of approximately 37 and results in the worst robustness across all metrics (AAAD, ASD, and ARIn). By contrast, perturbations to the mantissa bits have negligible impact on AAAD. Attacks on bit index 6 yield slightly higher AAAD than the mantissa bits, through still minor compared to bit index 1.

%It is observed that the attack on mantissa have almost no effect on the observed absolute accuracy difference, with the value of absolute accuracy difference observed is almost 0. Bit index 3 also observed to have a similar impact on both absolute accuracy difference and softmax difference, compared with the mantissa bits. 



%By contrast, perturbations to the mantissa bits have negligible impact on AAAD, with values close to zero, and produce robustness metrics nearly identical to those from bit index 3 in the exponent. Attacks on bit index 6 also yield slightly higher AAAD than the mantissa bits, through still minor compared to bit index 1.

%Based on Table \ref{table:robbothbitloc}, which is summarized in Figure \ref{fig:logabsdiff2}, we can see that a correlation between log absolute difference with average absolute accuracy difference and average softmax difference cannot be established. SEU in bit index 3 and 6 yield higher average accuracy difference compared to other bit index, excluding bit index 1. It is also observed that attack in bit index 0 results in a higher average softmax difference.

Notably, no clear correlation can be established between the log absolute difference and robustness metrics such as AAAD or ASD. Bit index 0 produces a moderate log absolute difference yet yields slightly higher ASD compared to mantissa flips. 

%Overall, these results confirm that SEUs affecting the exponent, particularly at bit index 1, are the most harmful to model robustness.

\subsection{Robustness by original bit value}

%Summary of absolute accuracy difference and softmax difference by original bit condition for Deterministic and Bayesian CNN
\begin{table}[H]
\centering
\begin{tabular}{l|rr|rr|rr}
\toprule
 & \multicolumn{2}{c|}{AAAD} & \multicolumn{2}{c|}{ASD} & \multicolumn{2}{c}{ARIn} \\
 & DNN & BNN & DNN & BNN & DNN & BNN \\
\midrule
0 & \B{0}{.0275} & \B{0}{.0455} & \B{0}{.2767} & \B{0}{.1734} & \B{0}{.1966} & \B{0}{.1268} \\
1 & \B{1}{.0002} & \B{1}{.0069} & \B{1}{.2612} & \B{1}{.0529} & \B{1}{.1847} & \B{1}{.0378} \\
\bottomrule
\end{tabular}
\caption{Robustness of Deterministic and Bayesian CNN model by original bit attacked}
\label{table:robbothorbit}
\end{table}

In this subsection, we also compare how the robustness is affected by the original bit value/state before the SEU. Table \ref{table:robbothorbit} shows that bitflip from 1 to 0 exhibits an extremely more robust model compared to bitflip from 0 to 1, especially on the absolute accuracy difference. It's not aligned with the claim from \cite{Hanif2021} that stated the opposite. 

This difference can be explained by the distribution of weights and biases in our models, which are typically close to 0 or 1. Consequently, the original bit value/state play an important role in determining the impact of SEUs. After a thorough observation, we saw that in all SEU injections involving bit index 1, the original bit value/state was always 0. 

%Since Subsection \ref{robattackedbit} already established that bit index 1 is highly vulnerable to SEUs, the observed lack of robustness in bitflip from 0 to 1 are consistent with these findings.

\section{Conclusion}
\label{sec:conclusion}


%\section{\uppercase{Conclusion}}

%Our study explores the double descent phenomenon in three recent deep learning models. QNNs exhibited clear double descent on CIFAR-10, while LTCs only showed some evidence of it, requiring further validation. SNNs did not display double descent. We found that learning rate schedulers, optimizers, and training epochs significantly influence double descent. Future work will involve broader experiments with varied parameters, optimizers, schedulers, and models, including FNNs, RNNs, and GANs.

%We believe that double descent reflects the behavior of dynamically learning features by the model during training, signifying that the model first learns shallow features and subsequently captures deeper, more complex features. Models exhibiting this pattern often demonstrate strong generalization. However, the absence of the double descent phenomenon does not necessarily indicate good or poor generalization performance. For instance, a model may quickly learn sufficient features, causing the test error curve to decline rapidly and stabilize at a low value without exhibiting a second increase. Conversely, the test error curve may take on a U-shaped pattern, as observed in SNN models trained without a learning rate scheduler. The reasons behind the lack of double descent phenomenon in SNNs, however, require further investigation.

%In summary, we believe that double descent is generally a favorable phenomenon for generalization. However, researchers should not explicitly aim to achieve double descent; instead, they should remain focused on the ultimate goal of machine learning—achieving better generalization. To this end, selecting appropriate hyperparameters and adopting efficient learning rate schedulers, as demonstrated in this study, can be effective strategies.


This research explored the performance of Bayesian Neural Networks (BNNs) robustness against Single Event Upsets (SEUs). Our initial findings show that the Deterministic CNN model is better in test accuracy compared to the Bayesian CNN model. On the other hand, Bayesian CNN model is more robust, looking at its performance in ARIn and Softmax Difference. Based on our experiment result, model with dropout is the most robust for Bayesian CNN. The Gaussian prior is the most robust. SEU against the prior center parameter is more robust than the spread parameter. Sigmoid offers better robustness in ASD and ARIn, RWG is the most robust in AAAD. Trade-off between test accuracy and robustness is observed. Convolutional layers are more robust than the fully connected layer. Attack on the first exponent bit is the most devastating; mantissa bit flips had a minimum effect. Flipping a bit from 1 to 0 is found to be far more robust.

%Future research could move beyond evaluating performance on unseen datasets (e.g., accuracy, F1 score) and instead analyze how SEUs affect the magnitudes and distributions of weights and biases in subsequent layers. Such an analysis may help identify sensitive components within the network and inform targeted mitigation strategies.

%Beyond architectural modifications, future work could also investigate distributional adjustments at the predictive level to improve robustness. One promising direction is applying Laplace smoothing to the output probability distributions to stabilize predictions under SEU-induced perturbations. It would be interesting if we can find any correlation between the smoothing factor and SEU robustness.




\bibliographystyle{apalike}
{\small
% \bibliography{PutawaraRevalda_mscthesis}}
\bibliography{references}}

% \newpage
% \section*{\uppercase{Appendix}}

% % If any, the appendix should appear directly after the
% % references without numbering, and not on a new page. To do so please use the following command:
% % \textit{$\backslash$section*\{APPENDIX\}}
% \begin{table}[h]
%     \caption{LTC Result Summary}
%     \label{ltc:double-descent}
%     \centering
%     \small
%     \begin{tabular}{p{0.65cm} p{1.9cm} p{1.2cm} p{0.8cm}}
%         \toprule
%         Model & \makecell[l]{Default \\ Setting} & \makecell[l]{Changing \\ Setting} & Result \\
%         \midrule
%         \multirow{4}{*}{LTCs} & \makecell[l]{Fashion-MNIST \\ SGD \\ with scheduler \\ 0\% label noise} & 1 layer & \(\times\) \\
%         & & 3 layers & \(\times\) \\
%         & & 5 layers & possible \\
%         & & 10 layers & \(\times\) \\
%         \cmidrule{2-4}
%         & \makecell[l]{Adam \\ with scheduler \\ 0\% label noise} & 5 layers & \(\times\)\\
%         \bottomrule
%     \end{tabular}
% \end{table}

% \begin{table} [h]
%     \caption{QNN Result Summary}
%     \label{qnn:double-descent}
%     \centering
%     \small
%     \begin{tabular}{p{0.65cm} p{2cm} p{2cm} p{0.8cm}}
%         \toprule
%         Model & \makecell[l]{Default \\ Setting} & \makecell[l]{Changing \\ Setting} & Result \\
%         \midrule
%         \multirow{5}{*}{QNNs} & \makecell[l]{CIFAR-10 \\ 4 layers \\ Adam \\ without scheduler \\ 50 epochs} & 0\% label noise & \(\checkmark\) \\
%         & & 10\% label noise & \(\checkmark\) \\
%         & & 20\% label noise & \(\checkmark\) \\
%         \cmidrule{2-4}
%         & \makecell[l]{CIFAR-10 \\ 4 layers \\ Adam \\ without scheduler \\ 10\% label noise \\ 80 epochs} & & \(\checkmark\) \\
%         \cmidrule{2-4}
%         & \makecell[l]{CIFAR-10 \\ 4 layers \\ 10\% label noise \\ 20 epochs} & \makecell[l]{Adam\\ no scheduler} & \(\times\) \\
%         & & \makecell[l]{Adam \\ with scheduler} & \(\times\) \\
%         & & \makecell[l]{SGD \\ no scheduler} & \(\checkmark\) \\
%         & & \makecell[l]{SGD \\ with scheduler} & \(\times\) \\
%         \cmidrule{2-4}
%         & \makecell[l]{MNIST \\ Adam \\ without scheduler \\ 10\% label noise \\ 50 epochs} & 4 layers & \(\times\) \\
%         & & 2 layers & \(\checkmark\) \\
%         \bottomrule
%     \end{tabular}
% \end{table}

% \begin{table}[h]
%     \caption{SNN Result Summary}
%     \label{snn:double-descent}
%     \centering
%     \small
%     \begin{tabular}{p{0.65cm} p{2cm} p{2cm} p{0.8cm}}
%         \toprule
%         Model & \makecell[l]{Default \\ Setting} & \makecell[l]{Changing \\ Setting} & Result \\
%         \midrule
%         \multirow{4}{*}{SNN} & \makecell[l]{MNIST \\ SGD \\ 10\% label noise \\ 50 epochs} & 2 layers, without scheduler & \(\times\) \\
%         & & \makecell[l]{2 layers\\ with scheduler} & \(\times\) \\
%         & & \makecell[l]{4 layers\\ without scheduler} & \(\times\) \\
%         & & \makecell[l]{4 layers\\ with scheduler} & \(\times\) \\
%         \bottomrule
%     \end{tabular}
% \end{table}

\end{document}

